%%=============================================================================
%% Conclusie
%%=============================================================================

\chapter{Conclusie}%
\label{ch:conclusie}

Dit onderzoek vertrok van de vraag hoe automatische netwerkdetectie kan bijdragen aan het up-to-date houden van netwerkdocumentatie en aan het realiseren van een betrouwbare single source of truth (SSOT) op geselecteerde WiFi-/AP-sites bij Citymesh. De kern van de bijdrage ligt niet enkel in de evaluatie van NetBox Discovery als bestaande tool, maar in het ontwerp en de uitrol van een hybride discoverypipeline waarin NetBox Discovery met Diode de inventarislaag voedt en \texttt{NetworkTool}, een eigen ontwikkeling, de LLDP-gebaseerde kabel- en topologielaag en de RUCKUS vSZ-koppeling afdekt. Op basis van een literatuurstudie (Hoofdstuk~\ref{ch:stand-van-zaken}), een gestructureerde casestudy met vier sites (Hoofdstuk~\ref{ch:casestudy}) en de implementatie van die hybride pipeline (Hoofdstuk~\ref{ch:implementatie}) kan de onderzoeksvraag concreet beantwoord worden.

De casestudy laat zien dat automatische discovery binnen de Citymesh-context effectief leidt tot een betere NetBox-documentatie (synthese in tabel~\ref{tab:synthese}). Op de drie sites waar kabelrelaties als kerntaak in scope zaten (Ridefort, Koninklijke Stallingen en WZC Riethove), werden ze van geheel of grotendeels ontbrekend hersteld tot volledig of zo volledig als operationeel haalbaar was; de enige resterende lacune is herleidbaar tot \'e\'en fysiek offline switch op WZC Riethove en zal in een volgende discovery aangevuld worden. Op Deleers Go4tech, waar de baseline al consistent was, corrigeerde de discovery enkele switches met een foutief device type, wat aantoont dat ook recent ingevoerde documentatie nog menselijke fouten kan bevatten. Daarnaast biedt de RUCKUS vSZ-koppeling een herhaalbare bron voor de vendor-specifieke AP-velden die via SNMP/LLDP onvoldoende betrouwbaar beschikbaar zijn (zie \S\ref{sec:vsz-via-api}).

Tegelijk maakt deze studie ook de grenzen van de standaard discoverystack zichtbaar. De NetBox Discovery-keten (Orb agent + Diode + NetBox) levert vandaag geen automatische kabelobjecten op basis van LLDP-buren, en biedt geen ingebouwde RUCKUS vSZ-integratie. Voor een volledig topologisch SSOT-resultaat op WiFi-/AP-sites was de \texttt{NetworkTool}-laag dus geen cosmetische toevoeging, maar een noodzakelijke aanvulling. De hybride pipeline is daardoor het eigenlijke oplossingsobject van deze bachelorproef: ze toont aan dat NetBox Discovery met Diode pas zijn volledige meerwaarde realiseert in combinatie met procesmatige randvoorwaarden (cleanup vóór herdiscovery, periodieke herhaling, een centrale secret manager) en met aanvullende automatisering voor datatypes die de standaard keten niet afdekt.

Beknopt geformuleerd geven de onderzoeksvragen het volgende beeld. Wat het probleem betreft, manifesteert verouderde NetBox-documentatie zich in de praktijk vooral als ontbrekende of foutieve kabelrelaties en als afwijkingen in device type en naam na infrastructuurwijzigingen; de oorzaken liggen bij de manuele CSV-workflow en bij niet-systematische opvolging van vervangingen. Wat de oplossing betreft, levert NetBox Discovery met Diode goede inventaris-resultaten op, terwijl topologie pas volledig wordt na een aanvullende LLDP-stap via \texttt{NetworkTool}. De geobserveerde verbetering ten opzichte van de baseline is op de drie geëvalueerde indicatoren consistent (zie tabel~\ref{tab:synthese}), met als enige restrictie dat fysiek niet-bereikbare devices niet kunnen worden meegenomen. Een duurzame SSOT vereist tot slot dat discovery, cleanup en validatie procesmatig verankerd worden, eerder dan ad hoc te worden toegepast; de gedetailleerde aanbevelingen hierover staan in \S\ref{sec:aanbevelingen}.

Voor verder onderzoek liggen er drie pistes voor de hand. Ten eerste is het zinvol om de aanvullende \texttt{NetworkTool}-stap voor LLDP-kabelobjecten op te volgen en te herevalueren wanneer NetBox Labs deze stap native afdekt in een toekomstige Diode- of Orb-versie. Ten tweede zou een uitbreiding naar meer sites en vendormixen (bijvoorbeeld Cisco of Aruba) toelaten om de generaliseerbaarheid van de bevindingen verder te toetsen. Ten derde verdient de integratie met een centrale secret manager en een formele change-management-flow bijkomende uitwerking, zodat de uitrol naar tientallen sites operationeel beheersbaar blijft.
